This project set out to test whether compensation for wild boar damage
raises tolerance toward wild boar among residents of two counties bordering
Giant Panda National Park, and, if not, why not. Across 625 households, we
find the opposite of the policy assumption: compensation participation is
associated with \emph{lower}, not higher, tolerance ($\beta = -0.109$,
$p = .003$), even after adjusting for the strong confound between
compensation status and county of residence ($\chi^2 = 77.1$, $p < .001$).
This compensation paradox is not explained by the five mediating constructs
we measured -- compensation anticipation, policy support, tangible cost,
intangible cost, and intangible benefit account for only 18\% of the total
association (RQ1) -- and it does not depend on how costs and benefits
interact with compensation status, since none of the three tested
interaction terms is significant (RQ2). The paradox also survives extensive
robustness testing: it is negative in all 8{,}192 alternative model
specifications we estimated, and the underlying structural model is largely
invariant between compensated and uncompensated residents, with one
exception -- conflict experience with wildlife more broadly matters more for
compensated residents' tolerance than for uncompensated residents' (RQ3).

Throughout, intangible cost -- residents' general evaluative attitude and
emotional reaction toward wild boar -- is the strongest single predictor of
tolerance in every specification and every subgroup we tested, consistently
outweighing compensation status itself. This points toward a broader
conclusion than a narrow one about compensation policy: efforts to improve
tolerance that focus only on offsetting tangible economic loss are targeting
a secondary driver. The primary driver, in this data, is how residents feel
about wild boar and their own capacity to cope with living alongside them --
dimensions that a compensation payment alone does not address.

Two things follow for how these results should be used. First, they should
be read as a robust, well-tested \emph{association}, not as evidence that
compensation payments themselves cause lower tolerance; the cross-sectional
design cannot rule out that residents who were already less tolerant were
also more likely to apply. Second, resolving that ambiguity -- and
identifying what, beyond the five mediators tested here, actually connects
compensation experience to tolerance -- is squarely a question for the
longitudinal, before/after study design recommended in
Section~\ref{discussion:next-steps}, not for further re-analysis of the
present cross-sectional data.
